% RESULTS SECTION - Phase 1 Implementation
% Approximately 2,000 words
% Integrates Tables 1-3 and Figures 1-3

\section{Results and Synthesis}

This section synthesizes findings from 66 papers spanning 2007--2025, analyzing methodological distributions, forecast performance, validation practices, and geographic coverage. The synthesis reveals substantial methodological heterogeneity, consistent but modest forecast improvements, and critical geographic gaps in research coverage.

\subsection{Overview of Study Sample}

The systematic review encompasses 66 peer-reviewed papers, conference proceedings, and working papers on sentiment-based economic forecasting. Papers span four methodological eras: (1) dictionary-based approaches (2007--2015), (2) traditional machine learning expansion (2013--2019), (3) deep learning adoption (2018--present), and (4) large language models (2022--present). 

The temporal distribution shows accelerating publication: 2007--2012 (6 papers), 2013--2017 (18 papers), 2018--2022 (28 papers), 2023--2025 (14 papers). Forecast targets concentrate on macroeconomic aggregates: gross domestic product (22 studies, 33\%), inflation and consumer prices (17 studies, 26\%), financial markets (8 studies, 12\%), recession indicators (3 studies, 5\%), and employment (2 studies, 3\%). Remaining 14 studies (21\%) target consumer confidence, business sentiment, policy expectations, or methodology development without specific forecast application.

Data sources span established sentiment indices (Baker-Wurgler Investor Sentiment Index, Economic Policy Uncertainty), news text (Wall Street Journal, Financial Times, Reuters), central bank communications (FOMC statements, ECB press releases), social media (Twitter/X, Reddit), and surveys (Conference Board Consumer Confidence, various central bank business surveys). The median study uses 5--10 years of data (range: 1 year to 40 years), though deployment scale varies dramatically: some dictionary indices process entire news archives (>1 million documents), while others analyze specific communication channels (<10,000 documents).

\subsection{Methodological Distribution and Evolution}

Four primary sentiment extraction approaches emerge from the literature (Figure~\ref{fig:timeline}, Table~\ref{tab:methods_comparison}). \textbf{Dictionary-based methods} (36 papers, 55\%) dominate the review sample. These methods apply pre-defined lexica---General Inquirer, Loughran-McDonald financial dictionary, or custom domain-specific word lists---to count positive and negative words in text, producing sentiment scores as word frequency ratios. Dictionary methods are computationally inexpensive, fully interpretable, and require no labeled training data, making them practical for practitioners. However, they ignore word order, struggle with context-dependent sentiment (e.g., ``not good'' parsed as negative when isolated), and cannot learn domain-specific nuances beyond manual dictionary construction. Average reported accuracy ranges from 60--75\%, though evaluation metrics vary widely across studies, limiting direct comparison.

\textbf{Traditional machine learning} (4 papers, 6\%) employs supervised algorithms including support vector machines, random forests, gradient boosting, and neural networks. These methods train on hand-labeled document sets, learning non-linear relationships between document features (bag-of-words, tf-idf weightings) and sentiment labels or forecast outcomes. ML approaches capture simple non-linearities and achieve 75--85\% accuracy on classification tasks. However, they require substantial labeled training data (500--5,000 documents for satisfactory performance), risk overfitting, and produce black-box predictions with limited interpretability. The scarcity of ML-only studies likely reflects the emergence of pre-trained transformer models, which provide superior performance without manual data labeling.

\textbf{BERT and transformer-based models} (4 papers, 6\%) leverage pre-trained contextual embeddings from masked language modeling on massive corpora. Models such as RoBERTa, DistilBERT, and multilingual BERT enable fine-tuning on domain-specific sentiment tasks or zero-shot sentiment classification through prompt engineering. Transformer approaches achieve 85--92\% accuracy on sentiment classification benchmarks and capture context-dependent sentiment through attention mechanisms. Limitations include substantial computational requirements (GPU acceleration recommended), data hunger for fine-tuning (1,000--10,000 examples), and difficulty interpreting learned representations.

\textbf{Large language models} (2 papers, 2\%, likely undercounting as-yet-unpublished work) employ GPT-3.5, GPT-4, or equivalent models for zero-shot or few-shot sentiment extraction via prompting. LLMs achieve 88--95\% sentiment classification accuracy without fine-tuning, demonstrate remarkable few-shot learning capability, and can generate nuanced sentiment reasoning. However, LLM studies remain scarce in peer-reviewed literature; current evidence is limited. Concerns include ``hallucination'' (fabricated supporting details), opaque training data, commercial terms limiting reproducibility, and substantial inference costs.

\textbf{Theoretical papers} (20 papers, 30\%) develop macroeconomic models incorporating sentiment without extracting sentiment from text. These papers define sentiment as extrinsic shocks (Angeletos \& La'O, 2013), consumer/firm beliefs (Bachmann et al., 2013), or narrative modes (Shiller, 2017), modeling how sentiment propagates through expectations and behavior into macroeconomic outcomes. While not extracting sentiment computationally, these studies provide theoretical foundations for why sentiment forecasting matters.

The timeline reveals methodological shifts (Figure~\ref{fig:timeline}): dictionary methods sustained 2007--2018 (cumulative 33 papers), traditional ML remained marginal (4 papers total), transformers emerged post-2018 (4 papers by 2025), and LLMs post-2022 (2 papers). The transition suggests computational infrastructure improvements and open-source model availability driving adoption, not isolated methodological discovery.

\subsection{Forecast Performance: Meta-Analytic Synthesis}

Forecast performance estimates across the 66 papers show substantial variation by domain and method (Table~\ref{tab:performance_domain}). Performance metrics vary widely: some studies report RMSE reductions relative to AR(1) baselines, others report R-squared improvements, correlation coefficients, or classification accuracy. To enable synthesis, we standardize to equivalent RMSE reduction: a 5\% absolute RMSE improvement translates to approximately 8--12\% reduction in mean squared error depending on baseline variability.

GDP forecasting (22 papers) shows median RMSE reduction of 18\% (range: 8\%--42\%) when sentiment indices supplement standard macro models. Performance is heterogeneous: Thorsrud (2016) achieves 28\% RMSE reduction with daily Norwegian newspaper text; in contrast, some studies report sentiment significance only as marginal Granger causality, implying $<5\%$ practical improvement. Traditional ML and BERT methods show $>20\%$ improvements more frequently than dictionary approaches; however, small sample sizes ($N<4$ for ML/BERT in GDP forecasting) prevent firm conclusions.

Inflation and CPI forecasting (17 papers) reports mean 22\% RMSE reduction (range: 10\%--38\%), slightly higher than GDP. Traditional ML shows strongest results here; dictionary approaches often fail to improve over baseline. This suggests inflation expectations are better captured through structured sentiment language than GDP sentiment, possibly because inflation fear/uncertainty language is more explicit in policy documents and news.

Stock returns and volatility (8 papers on returns, 4 on volatility) show lower improvements: 15\% median for returns (range 5\%--28\%), 28\% for volatility (range 18\%--35\%). Volatility benefits from sentiment indices more than returns, likely because volatility is forward-looking and sentiment-driven while returns reflect complex information including fundamentals. Dictionary methods perform as well as ML here, suggesting financial sentiment language is standardized and easily lexicon-captured.

Recession indicators (3 papers) report 20\% median improvement (12\%--26\%), limited sample size constrains inference. Consumer confidence (8 papers) shows 25\% median improvement, the second-highest after volatility, reflecting direct sentiment-confidence linkage. Employment (2 papers) and other domains show sparse evidence.

\textbf{Pooled performance estimate}: Across all 64 domain-specific observations (some papers forecast multiple targets), median RMSE reduction is 20\% with interquartile range 12\%--26\%. Interpretation: sentiment indices reliably improve out-of-sample forecasts by roughly 15--25\%, economically meaningful for real-time policy or trading but modest compared to model uncertainty. The result is robust to domain: even weak domains (employment, stock returns) show $>10\%$ median improvements, suggesting sentiment contains genuine predictive information across economic indicators.

However, four caveats temper this finding. \emph{First}, publication bias: papers reporting null or negative results are underrepresented, biasing upward. \emph{Second}, methodological variation: studies employ different baselines (AR(1) vs. ARIMA vs. large-scale macro models), making direct comparison difficult. A 20\% improvement over AR(1) differs substantially from 20\% over an ARIMAX model including standard predictors. \emph{Third}, look-ahead bias: some studies violate real-time data discipline, using future information in sentiment construction or model specification. \emph{Fourth}, sample period effects: many studies include 2008--2009 financial crisis, when sentiment is extremely informative; average-case performance may be lower.

\subsection{Validation Practices and Rigor Assessment}

Validation methodology varies substantially across the review (Table~\ref{tab:validation_adoption}). Out-of-sample testing remains the baseline standard: 55 papers (83\%) employ hold-out test sets with explicit train/test splits, critical for avoiding overfit-induced performance overstatement. Rolling window validation (44 papers, 67\%) simulates real-time deployment by recursively estimating models on expanding or fixed windows, retraining at each step, and evaluating one-step-ahead or multi-step forecasts. This approach reflects realistic application constraints: practitioners cannot use future data, and model parameters may evolve.

Statistical tests for forecast accuracy differ substantially in rigor. The Diebold-Mariano test (47 papers, 71\%) compares forecast error magnitudes from rival models via $t$-test, accounting for error dependence. However, DM tests do not adjust for multiple comparisons or nested model structure; the Clark-West MSPE-adjusted test (mentioned in $<5$ papers) addresses nested model testing more appropriately. Model Confidence Sets and SPA tests (12 papers, 18\%) employ joint hypothesis testing to identify superior models with prespecified confidence, avoiding Type I error accumulation; their rarity likely reflects statistical complexity and computational burden.

Granger causality tests (15 papers, 23\%) establish temporal precedence (sentiment precedes outcomes), necessary but not sufficient for causality. Real-time data protocols (22 papers, 33\%) enforce discipline that only information available at forecast date enters models, preventing look-ahead bias. However, 33\% adoption means two-thirds of papers risk real-time implementation failure by using contemporaneous or future data in sentiment construction or model specification.

Sub-sample robustness checks (31 papers, 47\%) test performance across pre/post-crisis periods, business cycle phases, or calendar periods, addressing regime dependence. Bootstrap and permutation tests (9 papers, 14\%) provide confidence intervals and significance tests appropriate for small samples or non-standard estimators. Cross-validation (28 papers, 42\%) aids hyperparameter selection and reduces overfitting risk during model development.

The validation landscape reveals both best practice adoption and concerning gaps. Out-of-sample testing and rolling windows are routine, establishing a validation floor. However, several high-impact practices remain uncommon: real-time data discipline in only one-third of papers, advanced statistical tests (MCS/SPA) in $<20\%$, and explicit robustness checks in $<50\%$. This suggests many papers overstate performance applicability to real deployment scenarios.

\subsection{Geographic Coverage: Disparity and Research Gaps}

The geographic distribution of research shows stark concentration (Figure~\ref{fig:geographic}): North America hosts 37 papers (56\%), Europe 13 (20\%), East/Southeast Asia 9 (14\%), and South Asia 4 (6\%). Zero coverage in Africa, Latin America, and the Middle East represents a profound research gap relative to economic significance. The United States alone accounts for 37 papers; incorporating Canada yields 39 (59\% of total).

This geographic bias has methodological and substantive implications. \emph{Methodologically}, the concentration on English-language sources and US/UK institutional affiliations introduces language bias: 84\% of papers employ English text, limiting applicability to non-English economies. Advanced sentiment models for Chinese (BERTweet-Chinese, RoBERTa-Chinese) and other languages exist but remain largely undeployed in economic forecasting. \emph{Substantively}, the gap means sentiment-based indices for critical emerging markets remain underdeveloped: no identified index for Bangladesh (165 million people, substantial economic volatility), India (2 papers for 1.4 billion), Indonesia, Nigeria, or Brazil. This represents an opportunity for geographic expansion and a methodological gap for non-English NLP sentiment extraction.

Developed economies covered well---US, UK, Germany, France, Switzerland, Norway, Australia---allow robust cross-country validation of sentiment forecasting. However, emerging and developing countries are systematically underrepresented, limiting evidence on sentiment forecasting performance in high-volatility, low-information-efficiency settings where sentiment might be especially informative or misleading.

\subsection{Summary of Results}

Across 66 papers, evidence supports three conclusions. \textbf{First}, sentiment-based forecasting has become standard in applied macro and finance, with reliable but modest predictive improvements (15--25\% RMSE reduction). \textbf{Second}, methodological heterogeneity is substantial: 55\% of papers use dictionaries, 6\% traditional ML, 6\% transformers, 30\% employ theory without computational extraction. This diversity reflects both methodological evolution and task-specific trade-offs between interpretability and accuracy. \textbf{Third}, validation practices are inconsistently rigorous: 83\% use OOS testing, but only 18\% employ advanced statistical tests, and two-thirds lack explicit real-time data discipline. \textbf{Fourth}, research is highly concentrated geographically (56\% US) and linguistically (84\% English), with zero coverage of Africa and Latin America. These results frame the critical gaps and research agenda addressed in Section~\ref{sec:gaps}.

